\textbf{Phần 1. Trắc nghiệm lựa chọn}
\setcounterpageref{ex}{0}

\begin{ex} %[1D6V21]
Số gồm 3 chục và 5 đơn vị được viết là số nào?
\choice
{\True $35$}
{$34$}
{$36$}
{$33$}
\loigiai{
Số gồm 3 chục và 5 đơn vị viết là $35$.
\textbf{Đáp số:} $35$.\par
\textbf{Đáp án của câu hỏi là A}
}
\end{ex}

\begin{ex} %[1D6V21]
Số gồm 7 chục và 0 đơn vị được viết là số nào?
\choice
{$60$}
{$80$}
{\True $70$}
{$69$}
\loigiai{
Số gồm 7 chục và 0 đơn vị viết là $70$.
\textbf{Đáp số:} $70$.\par
\textbf{Đáp án của câu hỏi là C}
}
\end{ex}

\begin{ex} %[1D6V21]
Số 68 gồm mấy chục và 8 đơn vị?
\choice
{$5$}
{\True $6$}
{$7$}
{$4$}
\loigiai{
Số 68 có chữ số hàng chục là 6, nghĩa là gồm $6$ chục và 8 đơn vị.
\textbf{Đáp số:} $6$.\par
\textbf{Đáp án của câu hỏi là B}
}
\end{ex}

\begin{ex} %[1D6V21]
Số 94 gồm 9 chục và mấy đơn vị?
\choice
{$3$}
{$5$}
{$2$}
{\True $4$}
\loigiai{
Số 94 có chữ số hàng đơn vị là 4, nghĩa là gồm 9 chục và $4$ đơn vị.
\textbf{Đáp số:} $4$.\par
\textbf{Đáp án của câu hỏi là D}
}
\end{ex}

\begin{ex} %[1D6V21]
Số tròn chục bé nhất có hai chữ số là số nào?
\choice
{\True $10$}
{$0$}
{$20$}
{$9$}
\loigiai{
Các số tròn chục có hai chữ số là $10, 20, 30, \dots, 90$. Số bé nhất là số $10$.
\textbf{Đáp số:} $10$.\par
\textbf{Đáp án của câu hỏi là A}
}
\end{ex}

\begin{ex} %[1D6V21]
Số tròn chục lớn nhất có hai chữ số là số nào?
\choice
{$80$}
{$100$}
{\True $90$}
{$89$}
\loigiai{
Số tròn chục lớn nhất có hai chữ số là số $90$.
\textbf{Đáp số:} $90$.\par
\textbf{Đáp án của câu hỏi là C}
}
\end{ex}

\begin{ex} %[1D6V21]
Số lớn nhất có hai chữ số là số nào?
\choice
{$98$}
{\True $99$}
{$100$}
{$97$}
\loigiai{
Số lớn nhất có hai chữ số là số $99$.
\textbf{Đáp số:} $99$.\par
\textbf{Đáp án của câu hỏi là B}
}
\end{ex}

\begin{ex} %[1D6V21]
Số bé nhất có hai chữ số khác nhau là số nào?
\choice
{$0$}
{$20$}
{$9$}
{\True $10$}
\loigiai{
Số 10 có hàng chục là 1, hàng đơn vị là 0 khác 1 nên là số bé nhất có hai chữ số khác nhau.
\textbf{Đáp số:} $10$.\par
\textbf{Đáp án của câu hỏi là D}
}
\end{ex}

\begin{ex} %[1D6V21]
Số lớn nhất có hai chữ số khác nhau là số nào?
\choice
{\True $98$}
{$97$}
{$99$}
{$96$}
\loigiai{
Chữ số hàng chục lớn nhất là 9, chữ số hàng đơn vị lớn nhất khác 9 là 8. Vậy số đó là $98$.
\textbf{Đáp số:} $98$.\par
\textbf{Đáp án của câu hỏi là A}
}
\end{ex}

\begin{ex} %[1D6V21]
Có tất cả bao nhiêu số tròn chục có hai chữ số?
\choice
{$8$}
{$10$}
{\True $9$}
{$7$}
\loigiai{
Các số tròn chục có hai chữ số là: $10, 20, 30, 40, 50, 60, 70, 80, 90$. Có tất cả $9$ số.
\textbf{Đáp số:} $9$.\par
\textbf{Đáp án của câu hỏi là C}
}
\end{ex}

\begin{ex} %[1D6V21]
Có tất cả bao nhiêu số có hai chữ số giống nhau?
\choice
{$8$}
{\True $9$}
{$10$}
{$7$}
\loigiai{
Các số có hai chữ số giống nhau là: $11, 22, 33, 44, 55, 66, 77, 88, 99$. Có tất cả $9$ số.
\textbf{Đáp số:} $9$.\par
\textbf{Đáp án của câu hỏi là B}
}
\end{ex}

\begin{ex} %[1D6V21]
Tìm một số có hai chữ số, biết chữ số hàng chục là 4 và chữ số hàng đơn vị lớn hơn chữ số hàng chục 2 đơn vị.
\choice
{$45$}
{$47$}
{$44$}
{\True $46$}
\loigiai{
Chữ số hàng chục là 4. Chữ số hàng đơn vị là: $4 + 2 = 6$. Vậy số cần tìm là $46$.
\textbf{Đáp số:} $46$.\par
\textbf{Đáp án của câu hỏi là D}
}
\end{ex}

\textbf{Phần 2. Trắc nghiệm đúng sai}
\setcounterpageref{ex}{0}

\begin{ex} %[1D6V21]
Xét tính đúng hoặc sai của các khẳng định toán học sau đây:
\choiceTF
{\True Kết quả của bài toán «Tìm một số có hai chữ số, biết chữ số hàng đơn vị là 3 và chữ số hàng chục gấp đôi chữ số hàng đơn vị» là $63$.}
{\True Kết quả của bài toán «Từ hai chữ số 2 và 7, viết được bao nhiêu số có hai chữ số khác nhau» là $2$.}
{Kết quả của bài toán «Từ ba chữ số 1, 0, 5, viết được bao nhiêu số có hai chữ số khác nhau» là $5$.}
{Số đứng liền sau của số 49 là số $51$.}
\loigiai{
\begin{itemize}
    \item Ý a) Đúng: Chữ số hàng chục là: $3 + 3 = 6$. Vậy số cần tìm là $63$.
\textbf{Đáp số:} $63$.
    \item Ý b) Đúng: Các số có hai chữ số khác nhau lập từ 2 và 7 là: $27$ và $72$. Có tất cả $2$ số.
\textbf{Đáp số:} $2$.
    \item Ý c) Sai: Kết quả đúng là $4$, không phải $5$. Chữ số hàng chục có thể chọn là 1 hoặc 5 (không thể chọn 0):\n- Hàng chục là 1: lập được $10, 15$.\n- Hàng chục là 5: lập được $50, 51$.\nCó tất cả $4$ số.
\textbf{Đáp số:} $4$.
    \item Ý d) Sai: Kết quả đúng là $50$, không phải $51$. Số liền sau của 49 là: $49 + 1 = 50$.
\textbf{Đáp số:} $50$.
\end{itemize}
\begin{itemize}
    \item a) \textbf{Đúng}.
    \item b) \textbf{Đúng}.
    \item c) \textbf{Sai}.
    \item d) \textbf{Sai}.
\end{itemize}
}
\end{ex}

\begin{ex} %[1D6V21]
Xét tính đúng hoặc sai của các bài toán và phép tính sau đây:
\choiceTF
{Số đứng liền trước của số 70 là số $68$.}
{Kết quả của bài toán «Tìm số tròn chục đứng liền trước của số 40» là $31$.}
{\True Kết quả của bài toán «Tìm số tròn chục đứng liền sau của số 70» là $80$.}
{\True Kết quả của bài toán «Từ ba chữ số 1, 0, 5, viết được bao nhiêu số có hai chữ số khác nhau» là $4$.}
\loigiai{
\begin{itemize}
    \item Ý a) Sai: Kết quả đúng là $69$, không phải $68$. Số liền trước của 70 là: $70 - 1 = 69$.
\textbf{Đáp số:} $69$.
    \item Ý b) Sai: Kết quả đúng là $30$, không phải $31$. Các số tròn chục liên tiếp là $20, 30, 40, 50$. Số tròn chục liền trước 40 là $30$.
\textbf{Đáp số:} $30$.
    \item Ý c) Đúng: Số tròn chục liền sau của 70 là $80$.
\textbf{Đáp số:} $80$.
    \item Ý d) Đúng: Chữ số hàng chục có thể chọn là 1 hoặc 5 (không thể chọn 0):\n- Hàng chục là 1: lập được $10, 15$.\n- Hàng chục là 5: lập được $50, 51$.\nCó tất cả $4$ số.
\textbf{Đáp số:} $4$.
\end{itemize}
\begin{itemize}
    \item a) \textbf{Sai}.
    \item b) \textbf{Sai}.
    \item c) \textbf{Đúng}.
    \item d) \textbf{Đúng}.
\end{itemize}
}
\end{ex}

\textbf{Phần 3. Trả lời ngắn}
\setcounterpageref{ex}{0}

\begin{ex} %[1D6V21]
Một bó que tính có 1 chục que tính. Bạn An có 5 bó que tính như thế. Hỏi bạn An có tất cả bao nhiêu que tính?
\shortans{50}
\loigiai{
\textbf{Đáp số:} $50$\par
5 bó que tính, mỗi bó 1 chục que: có tất cả 5 chục que tính, tức là $50$ que tính.
\textbf{Đáp số:} $50$.
}
\end{ex}

\begin{ex} %[1D6V21]
Mẹ mua 4 vỉ trứng gà, mỗi vỉ có 10 quả trứng. Hỏi mẹ mua tất cả bao nhiêu quả trứng gà?
\shortans{40}
\loigiai{
\textbf{Đáp số:} $40$\par
4 vỉ trứng, mỗi vỉ 10 quả: mẹ mua 4 chục quả trứng, tức là $40$ quả trứng.
\textbf{Đáp số:} $40$.
}
\end{ex}

\begin{ex} %[1D6V21]
Số 85 có chữ số hàng chục hơn chữ số hàng đơn vị bao nhiêu đơn vị?
\shortans{3}
\loigiai{
\textbf{Đáp số:} $3$\par
Chữ số hàng chục là 8, chữ số hàng đơn vị là 5: $8 - 5 = 3$.
\textbf{Đáp số:} $3$.
}
\end{ex}

\begin{ex} %[1D6V21]
Tìm một số có hai chữ số biết tổng hai chữ số bằng 7 và chữ số hàng đơn vị bằng 0.
\shortans{70}
\loigiai{
\textbf{Đáp số:} $70$\par
Chữ số hàng chục là: $7 - 0 = 7$. Vậy số đó là $70$.
\textbf{Đáp số:} $70$.
}
\end{ex}

\begin{ex} %[1D6V21]
Có bao nhiêu số có hai chữ số mà chữ số hàng đơn vị là 5?
\shortans{9}
\loigiai{
\textbf{Đáp số:} $9$\par
Các số đó là: $15, 25, 35, 45, 55, 65, 75, 85, 95$. Có tất cả $9$ số.
\textbf{Đáp số:} $9$.
}
\end{ex}

\begin{ex} %[1D6V21]
Có bao nhiêu số có hai chữ số mà chữ số hàng chục là 8?
\shortans{10}
\loigiai{
\textbf{Đáp số:} $10$\par
Các số đó là: $80, 81, 82, 83, 84, 85, 86, 87, 88, 89$. Có tất cả $10$ số.
\textbf{Đáp số:} $10$.
}
\end{ex}
